\subsection{CAR algebra and quasifree states}
We introduce an algebra of canonical anti-commutation relations (CAR algebra, for short) with a general one-particle Hilbert space following \cite{Araki1970,Binnenhei1995}.
Another style of, but equivalent, definition of a CAR algebra, can be found in, e.g., \cite{BratteliRobinson1997} (see Remark~\ref{rem:equivalence_CAR} below for the equivalence).
Let~$\mcal{K}$ be a complex Hilbert space of infinite dimension and $\Gamma$ be an anti-unitary involution on $\mcal{K}$.
The algebra $\mcal{C}_{0}(\mcal{K},\Gamma)$ is a~$\ast$-algebra over $\mbb{C}$ generated by $B(f)$, $f\in\mcal{K}$ subject to the relations
\begin{alignat*}{3}
	&B(\alpha f+\beta g)=\alpha B(f)+\beta B(g),\qquad && f, g\in\mcal{K},\quad \alpha,\beta\in\mbb{C},& \\
	&B(f)^{\ast}= B(\Gamma f),\qquad && f\in\mcal{K},&\\
	&\{B(f)^{\ast}, B(g)\}=(f,g)_{\mcal{K}}, \qquad && f,g\in\mcal{K},&
\end{alignat*}
where $\{\cdot,\cdot\}$ is the anti-commutator; $\{a,b\}:=ab+ba$.
It is known that the algebra $\mcal{C}_{0}(\mcal{K},\Gamma)$ admits a unique C$^{\ast}$-norm $\|\cdot\|$. We denote the C$^{\ast}$-completion by $\mcal{C}(\mcal{K},\Gamma)=\overline{\mcal{C}_{0}(\mcal{K},\Gamma)}^{\|\cdot\|}$ and call it the (self-dual) CAR algebra with one-particle space $(\mcal{K},\Gamma)$. To those who are more familiar with regarding $\ell^{2}(\mfrak{X})$ as a one-particle Hilbert space, we emphasize that we adopt a~``doubled'' space as a one-particle Hilbert space.

For a general C$^{\ast}$-algebra $\mfrak{A}$, a state over it is, by definition, a linear functional $\varphi\colon \mfrak{A}\to\mbb{C}$ satisfying the conditions that
\begin{enumerate}\itemsep=0pt
\item[1)] 	for every $A\in\mfrak{A}$, $\varphi (A^{\ast}A)\ge 0$ holds,
\item[2)] 	it is normalized:
		\begin{equation*}
			\|\varphi\|:=\sup\big\{|\varphi (A)|\, \big|\, \|A\|=1\big\}=1.
		\end{equation*}
\end{enumerate}
Note that, from these properties, it can be deduced that $\varphi (1)=1$ (see, e.g., \cite[Chapter~I, Section~9]{Takesaki1979}). Since any positive element $B\in\mfrak{A}$ admits an expression $B=A^{\ast}A$ with some $A\in\mfrak{A}$, the first condition is equivalently stated that $\varphi (B)\ge 0$ for all positive elements $B\in \mfrak{A}$.

\begin{Definition}A state $\varphi$ over a CAR algebra $\mcal{C}(\mcal{K},\Gamma)$ is said to be quasi-free if the moments admit Wick's formula, i.e.,
\begin{gather*}%\label{eq:Wick-property1}
	\varphi (B(f_{1})\cdots B(f_{2n+1}) ) =0, \\
%\label{eq:Wick-property2}
	\varphi (B(f_{1})\cdots B(f_{2n}) ) =(-1)^{n(n-1)/2}\sum_{\sigma}\operatorname{sgn} \sigma \prod_{i=1}^{n}\varphi (B(f_{\sigma (i)})B(f_{\sigma (i+n)}) ),
\end{gather*}
where $\sigma$ runs over permutations in $\mfrak{S}_{2n}$ such that $\sigma(1)<\cdots<\sigma (n)$ and $\sigma (i)<\sigma (i+n)$, $i=1,\dots, n$.
\end{Definition}

It is immediate from the definition that, for a quasi-free state $\varphi$ over $\mcal{C}(\mcal{K},\Gamma)$, a $2n$-point correlation function is expressed in terms of a Pfaffian so that
\begin{equation*}
	\varphi (B(f_{1})\cdots B(f_{2n}) )=\pf \mbb{A}_{\varphi}(f_{1},\dots, f_{2n}),
\end{equation*}
where $\mbb{A}_{\varphi}(f_{1},\dots, f_{2n})$ is the unique anti-symmetric $(2n\times 2n)$-matrix defined by
\begin{equation*}
	\mbb{A}_{\varphi}(f_{1},\dots, f_{2n})_{i,j}=\varphi (B(f_{i})B(f_{j}) ),\qquad 1\le i <j \le 2n.
\end{equation*}

By definition, a quasi-free state over a CAR algebra is uniquely determined by the two-point function. In fact, we have the following.
\begin{Lemma}[\cite{Araki1970}]\label{lem:quasi-free_covariance}
The collection of quasi-free states over $\mcal{C}(\mcal{K},\Gamma)$ is in one-to-one correspondence with the collection of operators $\mcal{Q}(\mcal{K},\Gamma)$ defined in~\eqref{eq:collection_covariance_ops}, under which an operator $S\in\mcal{Q}(\mcal{K},\Gamma)$ corresponds to the quasi-free state~$\varphi_{S}$ defined by
\begin{equation*}%\label{eq:quasifree_covariance}
	\varphi_{S}\big(B(f)^{\ast}B(g)\big)=(f,Sg)_{\mcal{K}},\qquad f,g\in\mcal{K}.
\end{equation*}
\end{Lemma}
We will give a proof of Lemma~\ref{lem:quasi-free_covariance} in Section~\ref{sect:existence_PfPP} for readers' convenience.

As we have announced, we work on the case when $\mcal{K}=\ell^{2}(\mfrak{X})\oplus\ell^{2}(\mfrak{X})$ equipped with $\Gamma$ given by (\ref{eq:Gamma}), and will adopt this pair $(\mcal{K},\Gamma)$ in the sequel without any specification.
The associated CAR algebra $\mcal{C}(\mcal{K},\Gamma)$ is generated by $a_{x}=B((0,e_{x}))$ and $a_{x}^{\ast}=B((e_{x},0))$, $x\in\mfrak{X}$.
Notice that the notation is compatible with the $\ast$-involution since $\Gamma (e_{x},0)=(0,e_{x})$.
Then, from the definition of a quasi-free state, we have
\begin{equation}
\label{eq:quasifree_correlation_function}
	\varphi_{S}\big(a_{x_{1}}^{\ast} \cdots a_{x_{n}}^{\ast}a_{x_{n}}\cdots a_{x_{1}}\big)=\pf [\mbb{K}_{S}(x_{i},x_{j}) ]_{1\le i,j\le n},
\end{equation}
where the matrix-valued function $\mbb{K}_{S}(\cdot,\cdot)\colon \mfrak{X}\times\mfrak{X}\to M(2;\mbb{C})$ was defined in~(\ref{eq:matrix_kernel_S}) associated with the operator~$S$. In fact, the values of the matrix-valued function $\mbb{K}_{S}(\cdot,\cdot)$ is written in terms of the quasi-free state as
\begin{equation*}
	\mbb{K}_{S}(x,y)=\left(
	\begin{matrix}
	\varphi_{S}(a_{x}^{\ast}a_{y}^{\ast}) & \varphi_{S}(a_{x}^{\ast}a_{y}) \\
	\varphi_{S}(a_{x}a_{y}^{\ast})-\delta_{x,y} & \varphi_{S}(a_{x}a_{y})
	\end{matrix}
	\right),\qquad x,y\in\mfrak{X}.
\end{equation*}
At this stage, we can find an analogy between this expectation value and a correlation function of a PfPP.

\begin{Remark}\label{rem:equivalence_CAR}
We may define functionals $a^{\ast}\colon \ell^{2}(\mfrak{X})\to \mcal{C}(\mcal{K},\Gamma)$ and $a\colon \ell^{2}(\mfrak{X})\to \mcal{C}(\mcal{K},\Gamma)$ by extending the assignments $e_{x}\mapsto a_{x}^{\ast}$ and $e_{x}\mapsto a_{x}$, $x\in\mfrak{X}$ linearly and anti-linearly, respectively. In other words, these functionals are defined by $a^{\ast}(u)=B((u,0))$, $a(u)=B((0,Ju))$, $u\in \ell^{2}(\mfrak{X})$.
Then, we find anti-commutation relations that are widely accepted~\cite{BratteliRobinson1997}:
\begin{gather}\label{eq:anti-commutation_BratteliRobinson}
	 \{a^{\ast}(u),a^{\ast}(v)\}=\{a(u),a(v)\}=0, \qquad
	 \{a(u),a^{\ast}(v)\}=(u,v)_{\ell^{2}(\mfrak{X})}, \qquad u,v\in \ell^{2}(\mfrak{X}).
\end{gather}
Conversely, when we have linear and anti-linear functional $a^{\ast}$, $a$ satisfying the anti-commutation relations (\ref{eq:anti-commutation_BratteliRobinson}), the generators of $\mcal{C}(\mcal{K},\Gamma)$ are recovered as $B((u,v))=a^{\ast}(u)+a(Jv)$, $(u,v)\in\mcal{K}$.
\end{Remark}

Let us consider a commutative algebra topologically generated by $a_{x}^{\ast}a_{x}$, $x\in\mfrak{X}$, which is identified with the algebra $C(\Omega)$ of continuous functions on $\Omega$ by the correspondence
\begin{equation*}
	\prod_{i=1}^{n}a_{x_{i}}^{\ast}a_{x_{i}}\mapsto \chi_{\Omega_{x_{1},\dots, x_{n}}},\qquad x_{1},\dots, x_{n}\in\mfrak{X}\colon \ \mbox{distinct.}
\end{equation*}
In the sequel, we regard $C(\Omega)$ as a subalgebra of $\mcal{C}(\mcal{K},\Gamma)$ under this correspondence.

Our strategy to prove Proposition~\ref{prop:existence_ppp} is to identify the above expectation value~(\ref{eq:quasifree_correlation_function}) with the correlation function of the desired PfPP, expecting that for a quasi-free state $\varphi_{S}$, there exists a probability measure $M^{S}$ on $(\Omega,\Sigma)$ so that the restriction of $\varphi_{S}$ on the subalgebra $C(\Omega)$ is identical to the integration with respect to $M^{S}$:
\begin{equation*}
	\varphi_{S}(F)=\int_{\Omega}F(\omega)M^{S}({\rm d}\omega),\qquad F\in C(\Omega).
\end{equation*}
Then, the probability measure $M^{S}$ is automatically a PfPP.
We will see in Section~\ref{sect:existence_PfPP} that this indeed happens to prove Proposition \ref{prop:existence_ppp}.

